On dispose $2n$ tas de pierres sur un cercle. Leurs tailles sont
\[
a_1, a_2, \ldots, a_{2n}.
\]
À chaque opération, on choisit deux tas voisins et on ajoute une pierre à
chacun d'eux.
\begin{enumerate}
  \item Montrer que la somme alternée
  \[
  S = a_1 - a_2 + a_3 - a_4 + \cdots + a_{2n-1} - a_{2n}
  \]
  est invariante.
  \item On part de $(1, 0, 0, \ldots, 0)$.

  Peut-on obtenir une configuration dans laquelle tous les tas ont la même
  taille ?
  \item Même question si l'on part de $(1, 1, 0, \ldots, 0)$.

  L'invariant suffit-il à conclure ?
\end{enumerate}
